% MATH3701 Supplementary Proofs — shared preamble
% % MATH3701 Supplementary Proofs — shared preamble
% % MATH3701 Supplementary Proofs — shared preamble
% \input{proof_preamble} at the top of each proof file

\documentclass[11pt,a4paper]{article}
\usepackage{amsmath,amssymb,amsthm}
\usepackage{mathtools}
\usepackage{bm}
\usepackage{geometry}
\geometry{margin=2.5cm}
\usepackage{hyperref}
\hypersetup{colorlinks=true, linkcolor=blue, urlcolor=blue}

% ---- Math shortcuts (matching slides/preamble.tex) ----
\newcommand{\E}{\mathbb{E}}
\newcommand{\Var}{\mathrm{Var}}
\newcommand{\Cov}{\mathrm{Cov}}
\newcommand{\Prob}{\mathbb{P}}
\newcommand{\R}{\mathbb{R}}
\newcommand{\N}{\mathcal{N}}
\newcommand{\bX}{\mathbf{X}}
\newcommand{\bx}{\mathbf{x}}
\newcommand{\by}{\mathbf{y}}
\newcommand{\bY}{\mathbf{Y}}
\newcommand{\bb}{\mathbf{b}}
\newcommand{\ba}{\mathbf{a}}
\newcommand{\bbeta}{\boldsymbol{\beta}}
\newcommand{\beps}{\boldsymbol{\epsilon}}
\newcommand{\bmu}{\boldsymbol{\mu}}
\newcommand{\btheta}{\boldsymbol{\theta}}
\newcommand{\bU}{\mathbf{U}}
\newcommand{\bJ}{\mathbf{J}}
\newcommand{\bI}{\mathbf{I}}
\newcommand{\bH}{\mathbf{H}}
\newcommand{\bK}{\mathbf{K}}
\newcommand{\bL}{\mathbf{L}}
\newcommand{\bM}{\mathbf{M}}
\newcommand{\bS}{\mathbf{S}}
\newcommand{\bW}{\mathbf{W}}
\newcommand{\bz}{\mathbf{z}}
\newcommand{\boldeta}{\boldsymbol{\eta}}
\newcommand{\logit}{\mathrm{logit}}
\newcommand{\iid}{\stackrel{\mathrm{iid}}{\sim}}
\newcommand{\tr}{\mathrm{tr}}

% ---- Theorem environments ----
\theoremstyle{plain}
\newtheorem{proposition}{Proposition}
\newtheorem{lemma}{Lemma}
\newtheorem{theorem}{Theorem}
\theoremstyle{definition}
\newtheorem{remark}{Remark}

% ---- Module branding ----
\newcommand{\moduleheader}{%
  \noindent\textbf{MATH3701 Statistical Modelling}\\
  School of Mathematics, University of Leeds\\[6pt]
}
 at the top of each proof file

\documentclass[11pt,a4paper]{article}
\usepackage{amsmath,amssymb,amsthm}
\usepackage{mathtools}
\usepackage{bm}
\usepackage{geometry}
\geometry{margin=2.5cm}
\usepackage{hyperref}
\hypersetup{colorlinks=true, linkcolor=blue, urlcolor=blue}

% ---- Math shortcuts (matching slides/preamble.tex) ----
\newcommand{\E}{\mathbb{E}}
\newcommand{\Var}{\mathrm{Var}}
\newcommand{\Cov}{\mathrm{Cov}}
\newcommand{\Prob}{\mathbb{P}}
\newcommand{\R}{\mathbb{R}}
\newcommand{\N}{\mathcal{N}}
\newcommand{\bX}{\mathbf{X}}
\newcommand{\bx}{\mathbf{x}}
\newcommand{\by}{\mathbf{y}}
\newcommand{\bY}{\mathbf{Y}}
\newcommand{\bb}{\mathbf{b}}
\newcommand{\ba}{\mathbf{a}}
\newcommand{\bbeta}{\boldsymbol{\beta}}
\newcommand{\beps}{\boldsymbol{\epsilon}}
\newcommand{\bmu}{\boldsymbol{\mu}}
\newcommand{\btheta}{\boldsymbol{\theta}}
\newcommand{\bU}{\mathbf{U}}
\newcommand{\bJ}{\mathbf{J}}
\newcommand{\bI}{\mathbf{I}}
\newcommand{\bH}{\mathbf{H}}
\newcommand{\bK}{\mathbf{K}}
\newcommand{\bL}{\mathbf{L}}
\newcommand{\bM}{\mathbf{M}}
\newcommand{\bS}{\mathbf{S}}
\newcommand{\bW}{\mathbf{W}}
\newcommand{\bz}{\mathbf{z}}
\newcommand{\boldeta}{\boldsymbol{\eta}}
\newcommand{\logit}{\mathrm{logit}}
\newcommand{\iid}{\stackrel{\mathrm{iid}}{\sim}}
\newcommand{\tr}{\mathrm{tr}}

% ---- Theorem environments ----
\theoremstyle{plain}
\newtheorem{proposition}{Proposition}
\newtheorem{lemma}{Lemma}
\newtheorem{theorem}{Theorem}
\theoremstyle{definition}
\newtheorem{remark}{Remark}

% ---- Module branding ----
\newcommand{\moduleheader}{%
  \noindent\textbf{MATH3701 Statistical Modelling}\\
  School of Mathematics, University of Leeds\\[6pt]
}
 at the top of each proof file

\documentclass[11pt,a4paper]{article}
\usepackage{amsmath,amssymb,amsthm}
\usepackage{mathtools}
\usepackage{bm}
\usepackage{geometry}
\geometry{margin=2.5cm}
\usepackage{hyperref}
\hypersetup{colorlinks=true, linkcolor=blue, urlcolor=blue}

% ---- Math shortcuts (matching slides/preamble.tex) ----
\newcommand{\E}{\mathbb{E}}
\newcommand{\Var}{\mathrm{Var}}
\newcommand{\Cov}{\mathrm{Cov}}
\newcommand{\Prob}{\mathbb{P}}
\newcommand{\R}{\mathbb{R}}
\newcommand{\N}{\mathcal{N}}
\newcommand{\bX}{\mathbf{X}}
\newcommand{\bx}{\mathbf{x}}
\newcommand{\by}{\mathbf{y}}
\newcommand{\bY}{\mathbf{Y}}
\newcommand{\bb}{\mathbf{b}}
\newcommand{\ba}{\mathbf{a}}
\newcommand{\bbeta}{\boldsymbol{\beta}}
\newcommand{\beps}{\boldsymbol{\epsilon}}
\newcommand{\bmu}{\boldsymbol{\mu}}
\newcommand{\btheta}{\boldsymbol{\theta}}
\newcommand{\bU}{\mathbf{U}}
\newcommand{\bJ}{\mathbf{J}}
\newcommand{\bI}{\mathbf{I}}
\newcommand{\bH}{\mathbf{H}}
\newcommand{\bK}{\mathbf{K}}
\newcommand{\bL}{\mathbf{L}}
\newcommand{\bM}{\mathbf{M}}
\newcommand{\bS}{\mathbf{S}}
\newcommand{\bW}{\mathbf{W}}
\newcommand{\bz}{\mathbf{z}}
\newcommand{\boldeta}{\boldsymbol{\eta}}
\newcommand{\logit}{\mathrm{logit}}
\newcommand{\iid}{\stackrel{\mathrm{iid}}{\sim}}
\newcommand{\tr}{\mathrm{tr}}

% ---- Theorem environments ----
\theoremstyle{plain}
\newtheorem{proposition}{Proposition}
\newtheorem{lemma}{Lemma}
\newtheorem{theorem}{Theorem}
\theoremstyle{definition}
\newtheorem{remark}{Remark}

% ---- Module branding ----
\newcommand{\moduleheader}{%
  \noindent\textbf{MATH3701 Statistical Modelling}\\
  School of Mathematics, University of Leeds\\[6pt]
}


\begin{document}

\moduleheader

\begin{center}
  {\Large\bfseries Supplementary Proof: The Smoothing Spline Solution}
\end{center}

\bigskip

\noindent This note proves two results used in Chapter~12 of the notes.
\begin{enumerate}
  \item \textbf{Proposition~12.1} (Roughness as a quadratic form): the
    roughness of a natural spline can be written as
    $J_\nu(f) = c_\nu\,\bb^T K_\nu \bb$ for explicit constants and
    matrices.
  \item \textbf{Proposition~12.2} (Smoothing spline solution): the
    minimiser of the penalised criterion $R_\nu(f,\lambda)$ satisfies a
    block linear system.
\end{enumerate}

\section*{Part 1: Roughness as a quadratic form}

\subsection*{Setup}

Let $x_1 < \cdots < x_n$ be $n$ distinct knots.  A natural cubic
spline ($\nu=2$) with knots at $x_1,\dots,x_n$ has the representation
\[
  f(t) = a_0 + a_1 t + \sum_{i=1}^n b_i\,|t-x_i|^3,
\]
with constraints $\sum b_i = 0$ and $\sum b_i x_i = 0$.  The roughness
is $J_2(f) = \int_{-\infty}^{\infty}[f''(t)]^2\,dt$.

\subsection*{Second derivative of the representation}

Differentiating twice,
\[
  \frac{d^2}{dt^2}|t-x_i|^3 = 6|t-x_i|.
\]
The polynomial part contributes nothing, so
\[
  f''(t) = 6\sum_{i=1}^n b_i\,|t-x_i|.
\]
The natural boundary conditions ($f''=0$ outside $[x_1,x_n]$) hold
automatically: for $t > x_n$, $f''(t) = 6(\sum b_i)\,t - 6\sum b_i x_i = 0$
by the two constraints.

\subsection*{Computing $J_2(f)$}

We use integration by parts to evaluate $J_2(f)$.  Differentiating
once more,
\[
  f'''(t) = 6\sum_{i=1}^n b_i\,\mathrm{sgn}(t-x_i)
\]
(a step function with jumps $12 b_i$ at each $x_i$), and in the
distributional sense,
\[
  f''''(t) = 12\sum_{i=1}^n b_i\,\delta(t-x_i).
\]
Since $f'' = 0$ outside $[x_1,x_n]$, two integrations by parts give
\begin{align*}
  J_2(f) = \int_{-\infty}^{\infty} [f''(t)]^2\,dt
  &= \int_{-\infty}^{\infty} f''(t)\cdot f''(t)\,dt \\
  &= \bigl[f''(t) f'(t)\bigr]_{-\infty}^{\infty}
     - \int_{-\infty}^{\infty} f'''(t)\,f'(t)\,dt \\
  &= 0 - \bigl[f'''(t) f(t)\bigr]_{-\infty}^{\infty}
     + \int_{-\infty}^{\infty} f''''(t)\,f(t)\,dt.
\end{align*}
The boundary terms vanish ($f''=0$ and $f$ is linear outside the knot
range, so $f'''=0$ there).  Substituting $f''''$,
\[
  J_2(f) = \int_{-\infty}^{\infty} 12\sum_{i=1}^n b_i\,\delta(t-x_i)\cdot f(t)\,dt
  = 12\sum_{i=1}^n b_i\,f(x_i).
\]
Now expand $f(x_i) = a_0 + a_1 x_i + \sum_k b_k |x_i-x_k|^3$.
Using $\sum b_i = 0$ and $\sum b_i x_i = 0$,
\[
  \sum_{i=1}^n b_i f(x_i)
  = a_0\underbrace{\sum b_i}_{0}
  + a_1\underbrace{\sum b_i x_i}_{0}
  + \sum_{i=1}^n\sum_{k=1}^n b_i b_k\,|x_i-x_k|^3.
\]
Therefore
\[
  J_2(f) = 12\sum_{i=1}^n\sum_{k=1}^n b_i b_k\,|x_i-x_k|^3
  = c_2\,\bb^T K_2\,\bb,
\]
where $c_2 = 12$ and $[K_2]_{ik} = |x_i-x_k|^3$. \hfill$\square$

\begin{remark}
  The analogous result for the linear case ($\nu=1$, $p=1$) is
  $J_1(f) = c_1\,\bb^T K_1\,\bb$ with $c_1=-2$ and
  $[K_1]_{ik}=|x_i-x_k|$.  The derivation is similar, using a single
  integration by parts and the representation $f(t)=a_0+\sum b_i|t-x_i|$
  with $\sum b_i = 0$.
\end{remark}

\section*{Part 2: The smoothing spline solution}

\subsection*{The optimisation problem}

By Proposition~11.3 (optimality of natural interpolants), the minimiser
$\hat f$ of the penalised criterion
\[
  R_\nu(f,\lambda) = \|\by - \hat{\mathbf{f}}\|^2 + \lambda J_\nu(f)
\]
is itself a natural spline with knots at $x_1,\dots,x_n$.
Using Part~1, with $\hat{\mathbf{f}} = K_\nu\bb + L_\nu\ba$ (the
matrix representation from the notes), the problem becomes:
\begin{align}
  \text{minimise} \quad
  &Q(\ba,\bb) = (\by - K\bb - L\ba)^T(\by - K\bb - L\ba)
    + \lambda^* \bb^T K\bb \label{eq:obj}\\
  \text{subject to} \quad
  &L^T\bb = \mathbf{0}, \label{eq:con}
\end{align}
where $\lambda^* = c_\nu\lambda$, $K = K_\nu$, $L = L_\nu$.

\subsection*{Unconstrained optimisation via Lagrange multipliers}

Introducing a Lagrange multiplier vector $\bmu$ for the constraint
\eqref{eq:con}, we minimise the Lagrangian
\[
  \mathcal{L}(\ba,\bb,\bmu)
  = Q(\ba,\bb) - 2\bmu^T L^T\bb.
\]
Setting partial derivatives to zero:
\begin{align}
  \frac{\partial\mathcal{L}}{\partial\ba}
  &= -2L^T(\by - K\bb - L\ba) = \mathbf{0}
  \implies L^T\by = L^T K\bb + L^T L\,\ba, \label{eq:d_a}\\[4pt]
  \frac{\partial\mathcal{L}}{\partial\bb}
  &= -2K^T(\by - K\bb - L\ba) + 2\lambda^* K\bb - 2L\bmu = \mathbf{0}.
  \label{eq:d_b}
\end{align}
Using the symmetry $K = K^T$, equation~\eqref{eq:d_b} simplifies to
\[
  K\by = K(K + \lambda^* I_n)\bb + KL\,\ba + L\bmu.
\]

\subsection*{Eliminating the Lagrange multiplier}

The constraint \eqref{eq:con} ($L^T\bb=\mathbf{0}$) means that
$L^T(K\bb)=0$, so pre-multiplying \eqref{eq:d_b} by $L^T$ gives
$L^T\by = L^T L\,\ba$, which is consistent with \eqref{eq:d_a} when
$K\bb$ is orthogonal to the column space of $L$.  Working through the
algebra, the system \eqref{eq:d_a}--\eqref{eq:d_b} can be assembled
into the single block linear system:
\begin{equation}
  \underbrace{
  \begin{bmatrix}
    0 & L^T \\
    L & K + \lambda^* I_n
  \end{bmatrix}
  }_{M_\lambda}
  \begin{bmatrix}
    \hat\ba \\ \hat\bb
  \end{bmatrix}
  =
  \begin{bmatrix}
    \mathbf{0} \\ \by
  \end{bmatrix},
  \label{eq:block}
\end{equation}
where the top block encodes the constraint $L^T\hat\bb = \mathbf{0}$
(since the equation $0\cdot\hat\ba + L^T\hat\bb = \mathbf{0}$ is
exactly \eqref{eq:con}).  The matrix $M_\lambda$ is symmetric and, for
$\lambda>0$, is invertible.  Hence
\[
  \begin{bmatrix}
    \hat\ba \\ \hat\bb
  \end{bmatrix}
  = M_\lambda^{-1}
  \begin{bmatrix}
    \mathbf{0} \\ \by
  \end{bmatrix}.
  \hfill\square
\]

\begin{remark}
  The system~\eqref{eq:block} is the same as equation~(12.11) in the
  notes.  In practice, \texttt{smooth.spline} in \texttt{R} solves an
  equivalent system using numerically stable algorithms that exploit
  the banded structure of $K$ after an appropriate change of basis.
\end{remark}

\end{document}
